\documentclass[12pt, a4paper]{article}

% ── PACKAGES ──────────────────────────────────────────────────
\usepackage[utf8]{inputenc}
\usepackage[T1]{fontenc}
\usepackage{geometry}
\geometry{margin=2.5cm}
\usepackage{hyperref}
\hypersetup{
    colorlinks=true,
    linkcolor=blue,
    urlcolor=blue,
    citecolor=blue,
    pdftitle={LumB DNMT3A/HDAC2 Co-Expression Coupling:
              A Subtype-Specific Chromatin Co-Repressor
              Circuit That Locks ER Output and Determines
              HDACi Sensitivity},
    pdfauthor={Eric Robert Lawson},
    pdfkeywords={DNMT3A, HDAC2, Luminal B, LumB, LumA,
                 co-repressor, chromatin lock, ER output,
                 ESR1, TFF1, HDACi, entinostat, breast cancer,
                 co-expression coupling, Pearson r,
                 METABRIC, GSE176078, attractor geometry,
                 Waddington landscape, OrganismCore,
                 endocrine resistance, epigenetic brake}
}
\usepackage{amsmath}
\usepackage{booktabs}
\usepackage{longtable}
\usepackage{array}
\usepackage{parskip}
\usepackage{microtype}
\usepackage{authblk}
\usepackage{titlesec}
\usepackage{fancyhdr}
\usepackage{xcolor}
\usepackage{float}
\usepackage{enumitem}
\usepackage{setspace}

% ── COLOURS ───────────────────────────────────────────────────
\definecolor{repobox}{RGB}{240,247,255}
\definecolor{findingbox}{RGB}{245,255,245}
\definecolor{convergentbox}{RGB}{255,250,235}
\definecolor{mechanismbox}{RGB}{250,245,255}
\definecolor{novelbox}{RGB}{255,248,220}
\definecolor{actionbox}{RGB}{230,255,230}
\definecolor{statbox}{RGB}{240,240,255}

% ── HEADER / FOOTER ───────────────────────────────────────────
\pagestyle{fancy}
\fancyhf{}
\lhead{\footnotesize CS-LIT-8: LumB DNMT3A/HDAC2
       Co-Expression Coupling}
\rhead{\footnotesize OrganismCore \textbar{} 2026-03-06}
\rfoot{\small Page \thepage}
\renewcommand{\headrulewidth}{0.4pt}
\setlength{\headheight}{24pt}
\addtolength{\topmargin}{-10pt}

% ── SECTION FORMATTING ────────────────────────────────────────
\titleformat{\section}
  {\large\bfseries}{\thesection}{1em}{}[\titlerule]
\titleformat{\subsection}
  {\normalsize\bfseries}{\thesubsection}{1em}{}
\titleformat{\subsubsection}
  {\small\bfseries}{\thesubsubsection}{1em}{}

% ══════════════════════════════════════════════════════════════
% TITLE
% ══════════════════════════════════════════════════════════════
\title{
    \vspace{-1em}
    \textbf{LumB DNMT3A/HDAC2 Co-Expression Coupling:}\\[0.5em]
    \large A Subtype-Specific Chromatin Co-Repressor
    Circuit That Locks ER Output in Luminal~B and
    Is Absent in Luminal~A --- Mechanistic Foundation
    for Differential HDAC Inhibitor Sensitivity\\[0.4em]
    \normalsize \textit{OrganismCore Framework ---
    Document CS-LIT-8}\\[0.3em]
    \small Derived from Waddington Attractor Geometry
    of 19,542 Single Cancer Cells (GSE176078)\\[0.3em]
    \small Replicated in METABRIC Bulk RNA-seq
    ($n=1{,}980$, $p=5.68\times10^{-56}$)
}

\author[1]{Eric Robert Lawson}
\affil[1]{
    OrganismCore\\
    \href{mailto:OrganismCore@proton.me}
         {OrganismCore@proton.me}\\
    ORCID:
    \href{https://orcid.org/0009-0002-0414-6544}
         {0009-0002-0414-6544}
}

\date{March 6, 2026}

% ══════════════════════════════════════════════════════════════
\begin{document}
% ══════════════════════════════════════════════════════════════

\maketitle
\thispagestyle{fancy}

% ── REPOSITORY POINTER ────────────────────────────────────────
\begin{center}
\colorbox{repobox}{
\begin{minipage}{0.88\textwidth}
\vspace{0.5em}
\centering
\textbf{Full computational record --- scripts, data
caches, reasoning artifacts, and raw logs:}\\[0.4em]
\href{https://github.com/Eric-Robert-Lawson/attractor-oncology}
{\texttt{https://github.com/Eric-Robert-Lawson/attractor-oncology}}
\\[0.3em]
\small Part of the OrganismCore BRCA Cross-Subtype
Analysis (BRCA-S8h, 2026-03-05).\\
Repository ID: 1172048282 \textbar{}
MIT License \textbar{}
Python \textbar{}
Public\\[0.2em]
All predictions locked from attractor geometry
\textbf{before} any literature review was performed.
\vspace{0.5em}
\end{minipage}}
\end{center}

\vspace{0.8em}

% ── PRIMARY FINDING BOX ───────────────────────────────────────
\begin{center}
\colorbox{statbox}{
\begin{minipage}{0.88\textwidth}
\vspace{0.5em}
\centering
\textbf{THE PRIMARY FINDING:}\\[0.5em]
DNMT3A and HDAC2 are co-expressed at
$\mathbf{r=+0.267}$ in Luminal~B breast cancer
vs $\mathbf{r=+0.071}$ in Luminal~A
($\mathbf{p=5.68\times10^{-56}}$, METABRIC,
$n=1{,}980$).\\[0.3em]
This $3.76\times$ difference in coupling strength
is the molecular signature of a chromatin
co-repressor circuit that is switched ON in
LumB and OFF in LumA. It suppresses ER
transcriptional output in LumB --- the
mechanistic origin of the LumA/LumB endocrine
resistance paradox and the mechanistic basis
for LumB-specific HDAC inhibitor sensitivity.
\vspace{0.5em}
\end{minipage}}
\end{center}

\vspace{0.8em}

% ══════════════════════════════════════════════════════════════
\begin{abstract}
% ══════════════════════════════════════════════════════════════
\noindent
Luminal~B (LumB) breast cancer is more resistant
to endocrine therapy and has worse prognosis than
Luminal~A (LumA) despite having higher ESR1
(oestrogen receptor) transcript. No published
biomarker framework has provided a mechanistic
causal explanation for this paradox at the
chromatin level. We report a LumB-specific
co-expression coupling between DNMT3A and HDAC2
that constitutes a chromatin co-repressor circuit
locking ER output specifically in LumB: Pearson
$r=+0.267$ in LumB versus $r=+0.071$ in LumA
($p=5.68\times10^{-56}$, METABRIC, $n=1{,}980$).

This coupling was identified by Waddington
attractor geometry applied to 19,542 single
cancer cells (GSE176078) and replicated in
METABRIC bulk RNA-seq ($n=1{,}980$). Its
downstream consequence is ER output decoupling:
TFF1, a canonical oestrogen-responsive gene, is
specifically and significantly lower in LumB
than LumA despite higher ESR1 transcript ---
the TFF1/ESR1 decoupling replicated in METABRIC
at $p=0.0019$ (CS-LIT-9/23,
\href{https://doi.org/10.5281/zenodo.18884234}
{DOI: 10.5281/zenodo.18884234}).

The mechanistic chain is: DNMT3A recruits HDAC2
to ER target gene promoters; HDAC2 deacetylates
histones H3K9ac and H3K27ac, compacting chromatin
and suppressing transcription; DNMT3A then
methylates the CpG islands at these promoters,
creating a durable silencing mark. The result is
an ER protein that is present and correctly
transcribed but unable to drive its downstream
target programme --- ER output is locked.

HDAC2 inhibition (entinostat, class~I HDACi)
directly disrupts the enzymatic component of
this circuit. By inhibiting HDAC2, entinostat
prevents chromatin compaction at ER target gene
promoters, restoring ER output in LumB
specifically. LumA patients do not have the
DNMT3A/HDAC2 co-repressor circuit; entinostat
has nothing to dissolve in LumA. This is the
mechanistic explanation for why entinostat
benefit is LumB-specific and why class-level
HR$+$ trials show modest effects (the LumA
patients dilute the signal).

Published literature establishes that DNMT3A
and HDAC interactions co-operate on co-repressor
complexes (Narcancer 2021), HDAC2 overexpression
contributes to endocrine resistance in breast
cancer (Springer 2024), and DNMT3A misregulation
associates with tamoxifen resistance (OAE 2024).
What is not published: the specific LumB vs LumA
co-expression coupling quantification
($r=+0.267$ vs $r=+0.071$,
$p=5.68\times10^{-56}$), the framing of this
as a subtype-specific chromatin lock that
explains the endocrine resistance paradox, and
the derived prediction that it determines HDACi
benefit stratification. Verdict: CONVERGENT-NOVEL.
\end{abstract}

\vspace{0.5em}
\noindent\textbf{Keywords:}
DNMT3A; HDAC2; Luminal B; LumB; LumA; co-repressor;
chromatin lock; ER output; ESR1; TFF1; HDACi;
entinostat; breast cancer; co-expression coupling;
Pearson~$r$; METABRIC; GSE176078; attractor geometry;
Waddington landscape; OrganismCore; endocrine
resistance; epigenetic brake; H3K9ac; H3K27ac

\newpage
\tableofcontents
\newpage

% ══════════════════════════════════════════════════════════════
\section{Document Metadata}
% ══════════════════════════════════════════════════════════════

\begin{center}
\begin{tabular}{ll}
\toprule
\textbf{Field} & \textbf{Value} \\
\midrule
Document ID & CS-LIT-8 \\
Parent document & BRCA-S8h (2026-03-05) \\
Verdict & CONVERGENT-NOVEL \\
Date & 2026-03-06 \\
Author & Eric Robert Lawson \\
ORCID & 0009-0002-0414-6544 \\
Organisation & OrganismCore \\
Contact & OrganismCore@proton.me \\
Repository &
\href{https://github.com/Eric-Robert-Lawson/attractor-oncology}
{attractor-oncology} (ID: 1172048282) \\
Derivation data & GSE176078 (19,542 single cancer cells) \\
Replication data & METABRIC ($n=1{,}980$) \\
Primary statistic &
$r=+0.267$ LumB vs $r=+0.071$ LumA,
$p=5.68\times10^{-56}$ \\
Downstream consequence &
TFF1/ESR1 decoupling $p=0.0019$ \\
CS-LIT-9/23 DOI &
\href{https://doi.org/10.5281/zenodo.18884234}
{10.5281/zenodo.18884234} \\
CS-LIT-2 DOI &
\href{https://doi.org/10.5281/zenodo.18884158}
{10.5281/zenodo.18884158} \\
CS-LIT-15 DOI &
\href{https://doi.org/10.5281/zenodo.18893147}
{10.5281/zenodo.18893147} \\
License & CC BY 4.0 \\
Biological contradictions & 0 \\
\bottomrule
\end{tabular}
\end{center}

% ══════════════════════════════════════════════════════════════
\section{Context: The LumA/LumB Endocrine
Resistance Paradox}
% ══════════════════════════════════════════════════════════════

\subsection{The clinical paradox}

Luminal~B breast cancer is classified as
ER-positive by IHC, shares the same receptor
classification as Luminal~A, yet:

\begin{itemize}[noitemsep]
    \item Has \textit{higher} ESR1 transcript
    than LumA (confirmed in GSE176078 and
    METABRIC)
    \item Has \textit{worse} endocrine therapy
    response than LumA
    \item Has \textit{worse} prognosis than LumA
    despite equivalent receptor status
    \item Has \textit{higher} Ki67 and
    proliferation rate
\end{itemize}

The standard clinical explanation for LumB
endocrine resistance invokes CDK4/6-driven
proliferation, PIK3CA mutations, and ESR1
mutations emerging on therapy. These explain
resistance at the signalling and cell cycle
level. They do not explain why LumB has more
ER transcript but less ER output.

\subsection{The attractor geometry diagnosis}

In the OrganismCore Waddington framework, LumB
is classified as a \textbf{Type~1 slope arrest
with epigenetic brake}: the cell has committed
to luminal identity (FOXA1 present, ESR1 present,
luminal programme intact) but a specific chromatin
remodelling event has applied a brake to the ER
output programme downstream of the ER protein
itself. The ER is present. The ER is not fully
reading out its target genes. The brake is at
the chromatin level, acting between the ER protein
and its downstream gene targets.

This framing generates a testable prediction:
there should be a measurable co-repressor
signature in LumB that is absent in LumA --- and
it should be specifically associated with ER
target gene suppression.

The DNMT3A/HDAC2 co-expression coupling is
that signature.

% ══════════════════════════════════════════════════════════════
\section{The Finding: DNMT3A/HDAC2
Co-Expression Coupling}
% ══════════════════════════════════════════════════════════════

\subsection{Primary statistics}

\begin{center}
\colorbox{findingbox}{
\begin{minipage}{0.88\textwidth}
\vspace{0.5em}
\centering
\textbf{DNMT3A/HDAC2 CO-EXPRESSION COUPLING}\\[0.5em]
\begin{tabular}{lcc}
\toprule
\textbf{Subtype} &
\textbf{Pearson $r$} &
\textbf{Dataset} \\
\midrule
Luminal~B & $+0.267$ & METABRIC ($n=1{,}980$) \\
Luminal~A & $+0.071$ & METABRIC ($n=1{,}980$) \\
\midrule
\multicolumn{3}{c}{Difference: $p=5.68\times10^{-56}$} \\
\multicolumn{3}{c}{Ratio of coupling: $3.76\times$} \\
\bottomrule
\end{tabular}
\vspace{0.5em}
\end{minipage}}
\end{center}

\vspace{0.5em}

\subsection{What the coupling means biologically}

A Pearson $r$ measures the linear co-expression
of two genes across a population of samples.
When DNMT3A and HDAC2 are co-expressed at
$r=+0.267$ in LumB:

\begin{itemize}
    \item Cells with higher DNMT3A also tend to
    have higher HDAC2 in the same sample
    \item This is the statistical signature of
    co-recruitment: both proteins are being
    upregulated together, consistent with being
    regulated by the same upstream signal or
    being directly co-assembled into a functional
    complex
    \item Co-recruited DNMT3A and HDAC2 at the
    same genomic loci constitutes a co-repressor
    complex --- the biochemical mechanism is
    established in published literature
    (Narcancer 2021)
\end{itemize}

When the same $r$ is only $+0.071$ in LumA:

\begin{itemize}
    \item DNMT3A and HDAC2 expression is
    uncorrelated in LumA
    \item They are not being co-recruited to
    the same targets
    \item The co-repressor circuit is not active
    in LumA
    \item This is not a continuous gradation ---
    the $p=5.68\times10^{-56}$ for the difference
    confirms these are statistically distinct
    biological states, not a spectrum
\end{itemize}

\subsection{The genomic targets of the complex}

The DNMT3A/HDAC2 co-repressor complex in LumB
is predicted to act at ER target gene promoters.
The evidence for this localisation:

\begin{enumerate}
    \item TFF1 (a canonical ER target gene) is
    specifically decoupled from ESR1 in LumB:
    higher ESR1, lower TFF1, lower TFF1/ESR1
    ratio ($p=0.0019$, METABRIC). The only
    locus at which this decoupling would appear
    is the TFF1 promoter --- downstream of the
    ER protein, at the chromatin level.

    \item HDAC2 activity depletes H3K9ac and
    H3K27ac at target promoters. These marks
    are required for ER-driven transcription.
    Depleting them suppresses ER target gene
    output without affecting ER transcript or
    protein levels.

    \item DNMT3A adds CpG methylation to the
    same deacetylated promoters, creating a
    durable silencing mark that persists through
    cell division --- converting a transient
    chromatin state into a stable epigenetic
    memory.
\end{enumerate}

This combination --- HDAC2 deacetylation + DNMT3A
methylation at ER target gene promoters --- is
the molecular definition of the chromatin lock.
It is present in LumB ($r=+0.267$). It is absent
in LumA ($r=+0.071$).

% ══════════════════════════════════════════════════════════════
\section{Downstream Consequence: TFF1/ESR1
Decoupling}
% ══════════════════════════════════════════════════════════════

The downstream consequence of the DNMT3A/HDAC2
lock is directly measurable as the TFF1/ESR1
ratio. This is the subject of CS-LIT-9/23
(\href{https://doi.org/10.5281/zenodo.18884234}
{DOI: 10.5281/zenodo.18884234}); it is
summarised here as the functional readout of
the chromatin circuit described in this document.

\begin{center}
\colorbox{mechanismbox}{
\begin{minipage}{0.88\textwidth}
\vspace{0.5em}
\textbf{The ER output chain in normal luminal
cells (LumA):}\\[0.3em]
ESR1 transcript $\rightarrow$ ER protein
$\rightarrow$ ER binds oestrogen $\rightarrow$
ER binds ERE at TFF1 promoter (with FOXA1
as pioneer) $\rightarrow$ TFF1 transcribed\\[0.5em]
\textbf{The ER output chain in LumB
(DNMT3A/HDAC2 lock active):}\\[0.3em]
ESR1 transcript $\uparrow$ (higher than LumA)
$\rightarrow$ ER protein present $\rightarrow$
ER binds oestrogen $\rightarrow$ ER attempts
to bind ERE at TFF1 promoter $\rightarrow$
\textbf{HDAC2 has deacetylated H3K9ac/H3K27ac}
$\rightarrow$ \textbf{DNMT3A has methylated
the CpG island} $\rightarrow$ chromatin is
compacted, inaccessible $\rightarrow$ TFF1
transcription suppressed\\[0.5em]
\textbf{Result:} ESR1 high, TFF1 low.
TFF1/ESR1 ratio is decoupled. The ER is
present and transcribed but locked out of
its own target genes.
\vspace{0.5em}
\end{minipage}}
\end{center}

\begin{center}
\begin{tabular}{p{3.5cm}p{4cm}p{4.5cm}}
\toprule
\textbf{Measurement} &
\textbf{Result} &
\textbf{Dataset} \\
\midrule
TFF1/ESR1 ratio
(LumB vs LumA) &
Significantly lower in LumB
despite higher ESR1 &
METABRIC $n=1{,}980$
$p=0.0019$ \\[0.5em]
DNMT3A/HDAC2
coupling &
$r=+0.267$ LumB vs
$r=+0.071$ LumA &
METABRIC $n=1{,}980$
$p=5.68\times10^{-56}$ \\[0.5em]
Discovery cohort &
Confirmed in scRNA-seq
(19,542 cells) &
GSE176078
(26 patients) \\
\bottomrule
\end{tabular}
\end{center}

These two findings --- the DNMT3A/HDAC2 coupling
(this document, CS-LIT-8) and the TFF1/ESR1
decoupling (CS-LIT-9/23) --- are the cause and
the consequence of the same chromatin event,
measured independently.

% ══════════════════════════════════════════════════════════════
\section{Therapeutic Prediction: HDACi Benefit
Is LumB-Specific}
% ═══════════════════════════════════════════��══════════════════

\subsection{The mechanism of entinostat in LumB}

Entinostat is a class~I HDAC inhibitor
(primarily HDAC1, HDAC2, HDAC3). It inhibits
HDAC2 directly. The prediction from the circuit:

\begin{center}
\colorbox{mechanismbox}{
\begin{minipage}{0.88\textwidth}
\vspace{0.5em}
\textbf{In LumB (lock present):}\\[0.3em]
Entinostat $\rightarrow$ HDAC2 inhibited
$\rightarrow$ H3K9ac and H3K27ac restored
at ER target gene promoters $\rightarrow$
chromatin opens $\rightarrow$ ER can now
drive TFF1, GATA3, and downstream
target transcription $\rightarrow$
ER output restored $\rightarrow$
endocrine therapy now fully engages
an active receptor\\[0.5em]
\textbf{In LumA (lock absent):}\\[0.3em]
Entinostat $\rightarrow$ HDAC2 inhibited
$\rightarrow$ there are no HDAC2-compacted
ER target gene promoters to reopen $\rightarrow$
ER output is already unobstructed $\rightarrow$
no mechanism for improving ET response $\rightarrow$
no benefit
\vspace{0.5em}
\end{minipage}}
\end{center}

\subsection{Why this explains the class-level
trial results}

The published literature on entinostat in
HR$+$ breast cancer:

\begin{center}
\begin{tabular}{p{3.5cm}p{4cm}p{4.5cm}}
\toprule
\textbf{Trial / Source} &
\textbf{Result} &
\textbf{Framework interpretation} \\
\midrule
E2112 (SWOG/NRG,
entinostat
$+$ exemestane
vs exemestane,
unselected HR$+$) &
Negative for OS and PFS &
Unselected HR$+$ $\approx$
50\% LumA + 50\% LumB.
LumA non-responders dilute
LumB signal to null. \\[0.5em]
Meta-analysis
May 2025 (Springer,
pooled 4 RCTs,
$n=1{,}371$) &
PFS HR$=0.80$, $p=0.01$.
No OS benefit. &
Pooling across trials:
overall HR$+$ population.
LumB signal partially
visible but diluted. \\[0.5em]
China NMPA approval
(April 2024) &
Approved in HR$+$/HER2$-$
advanced breast cancer &
Consistent with modest
population-level benefit
in an HR$+$ class that
includes LumB. \\
\bottomrule
\end{tabular}
\end{center}

\subsection{Testable prediction}

\begin{center}
\colorbox{actionbox}{
\begin{minipage}{0.88\textwidth}
\vspace{0.5em}
\textbf{The falsifiable prediction:}\\[0.4em]
In any entinostat trial in HR$+$ breast cancer,
the LumB subgroup PFS HR will be substantially
below the class-level HR, and the LumA subgroup
PFS HR will be approximately $1.0$.\\[0.4em]
\textbf{Immediate testing opportunity:}\\
NCT07235618 (entinostat $+$ fulvestrant,
post-CDK4/6i failure, Sun Yat-sen University,
$N=50$, start January 2026, PI: Dr.\ Shusen Wang).
See CS-LIT-15,
\href{https://doi.org/10.5281/zenodo.18893147}
{DOI: 10.5281/zenodo.18893147}.
\vspace{0.5em}
\end{minipage}}
\end{center}

% ══════════════════════════════════════════════════════════════
\section{What Is and Is Not in the Published
Literature}
% ══════════════════════════════════════════════════════════════

\subsection{What IS in the literature}

\begin{enumerate}
    \item \textbf{DNMT3A/HDAC co-repressor
    complex biochemistry (Narcancer 2021):}
    DNMT3A and HDAC activity co-operate on
    co-repressor complexes to silence target
    genes. The biochemical basis for the complex
    is published. The subtype-specific
    quantification in breast cancer is not.

    \item \textbf{HDAC2 in breast cancer endocrine
    resistance (Springer 2024):}
    HDAC2 overexpression is associated with
    endocrine resistance in breast cancer. The
    mechanism of resistance is at the chromatin
    level. The connection to a specific LumB
    co-expression circuit is not published.

    \item \textbf{DNMT3A and tamoxifen resistance
    (OAE Publishing 2024):}
    DNMT3A misregulation is associated with
    tamoxifen resistance in breast cancer. The
    connection to HDAC2 co-expression or the
    LumB-specific circuit is not published.
\end{enumerate}

\subsection{What is NOT in the literature}

\begin{center}
\colorbox{novelbox}{
\begin{minipage}{0.88\textwidth}
\vspace{0.5em}
\textbf{The following have no published equivalent
as of 2026-03-06:}
\begin{itemize}[noitemsep]
    \item The quantification of the DNMT3A/HDAC2
    Pearson $r$ within LumB ($+0.267$) and LumA
    ($+0.071$) separately in any breast cancer
    dataset.
    \item The $p=5.68\times10^{-56}$ for the
    subtype-specific difference.
    \item The framing of this as a binary circuit
    (ON in LumB, OFF in LumA) rather than a
    continuous gradation.
    \item The derivation that this circuit is the
    causal mechanism of TFF1/ESR1 decoupling
    (ER output suppression) in LumB.
    \item The prediction that entinostat benefit
    is LumB-specific because the HDAC2 component
    of this circuit is the drug's target.
    \item The circuit as the causal explanation
    for why LumB has more ESR1 transcript than
    LumA but less ER functional output.
\end{itemize}
\vspace{0.5em}
\end{minipage}}
\end{center}

% ══════════════════════════════════════════════════════════════
\section{Position in the Published Series}
% ══════════════════════════════════════════════════════════════

CS-LIT-8 is the mechanistic foundation document
for the LumB chromatin lock. Its position in the
evidence chain:

\begin{center}
\begin{tabular}{p{2.5cm}p{8.5cm}}
\toprule
\textbf{Document} & \textbf{Role} \\
\midrule
CS-LIT-8\\(this document) &
The chromatin circuit: DNMT3A/HDAC2
co-expression coupling in LumB
($r=+0.267$ vs $r=+0.071$, $p=5.68\times10^{-56}$).
The molecular cause. \\[0.5em]
CS-LIT-9/23\\
\href{https://doi.org/10.5281/zenodo.18884234}
{DOI: 10.5281/zenodo.18884234} &
The functional consequence: TFF1/ESR1 decoupling
in LumB ($p=0.0019$). The measurable downstream
output of the circuit. \\[0.5em]
CS-LIT-15\\
\href{https://doi.org/10.5281/zenodo.18893147}
{DOI: 10.5281/zenodo.18893147} &
The drug intervention: entinostat disrupts the
HDAC2 component of the circuit, restoring ER
output in LumB specifically. \\[0.5em]
CS-LIT-22\\
\href{https://doi.org/10.5281/zenodo.18892788}
{DOI: 10.5281/zenodo.18892788} &
The patient selection instrument: FOXA1/EZH2
dual IHC classifies LumA vs LumB at the
pathology bench level. \\[0.5em]
CS-LIT-2\\
\href{https://doi.org/10.5281/zenodo.18884158}
{DOI: 10.5281/zenodo.18884158} &
The classification framework: LumB is Lock
Type~1 with an epigenetic brake. CS-LIT-8
is the molecular evidence for the brake. \\
\bottomrule
\end{tabular}
\end{center}

% ══════════════════════════════════════════════════════════════
\section{Locked Statement}
% ══════════════════════════════════════════════════════════════

\begin{center}
\colorbox{findingbox}{
\begin{minipage}{0.88\textwidth}
\vspace{0.6em}
\textbf{Stated once, precisely, without
inflation:}\\[0.5em]
In Luminal~B breast cancer, DNMT3A and HDAC2
are co-expressed at Pearson $r=+0.267$, compared
to $r=+0.071$ in Luminal~A ($p=5.68\times10^{-56}$,
METABRIC, $n=1{,}980$). This $3.76\times$ difference
in coupling strength is the statistical signature
of a chromatin co-repressor circuit that is active
in LumB and inactive in LumA. The circuit acts at
ER target gene promoters: HDAC2 deacetylates
H3K9ac/H3K27ac, DNMT3A methylates CpG islands,
and together they compact chromatin and suppress
ER transcriptional output despite elevated ESR1
transcript --- producing the TFF1/ESR1 decoupling
confirmed in METABRIC ($p=0.0019$). The published
literature confirms that DNMT3A/HDAC co-repressor
biochemistry exists (Narcancer 2021), that HDAC2
drives endocrine resistance in breast cancer
(Springer 2024), and that DNMT3A misregulation
associates with tamoxifen resistance (OAE 2024).
What is not published is the subtype-specific
co-expression coupling quantification, the
binary circuit framing (ON in LumB, OFF in LumA),
and the derived prediction that entinostat benefit
is LumB-specific because the HDAC2 component of
this circuit is the drug's target. Verdict:
CONVERGENT-NOVEL.
\vspace{0.5em}
\end{minipage}}
\end{center}

\end{document}
